\section{A ground state at the minimum over all sectors}
\label{sec:sector}\label{sec:foundations}

The local comparisons in the proof may replace electrons in a region by a
different number of electrons.  We therefore first reduce to a ground state
whose energy is minimal among \emph{all} particle-number sectors.  A coarse
number bound makes this reduction possible; no convexity of the sector
energies is needed.

\subsection{The Coulomb form and strict binding}

Use the Hilbert spaces and form domains defined in the Introduction,
writing $n$ for a variable particle number.  The quadratic form associated
with \eqref{eq:hamiltonian} is
\begin{equation}\label{eq:sector-form}
 q_n[\psi]
 =\frac12\sum_{i=1}^n\|\nabla_i\psi\|_2^2
  -\sum_{i=1}^n\int V(x_i)|\psi|^2
  +\sum_{1\leq i<j\leq n}\int K(x_i-x_j)|\psi|^2.
\end{equation}
We also evaluate this expression on the ambient spin-valued $H^1$ space
when a localization preserves antisymmetry only within groups of variables.
Sector-energy bounds will be applied only to the groups that remain
antisymmetric.
Hardy's inequality in three dimensions, followed by Cauchy--Schwarz,
shows that each Coulomb term is infinitesimally form bounded relative to
the kinetic energy.  For a pair term, apply the same argument to the
relative coordinate $x_i-x_j$.  Thus, for each fixed system and sector,
the sum of the absolute Coulomb forms is bounded by
\[
 \eta\sum_i\|\nabla_i\psi\|_2^2+C_\eta\|\psi\|_2^2
 \qquad(\eta>0).
\]
Consequently \eqref{eq:sector-form} is closed and bounded below on
$\mathcal Q_n$.  Its associated self-adjoint operator is the Hamiltonian
$H_n$ in the Introduction, and
\[
 E_n=\inf_{\substack{\psi\in\mathcal Q_n\\\|\psi\|_2=1}}
             q_n[\psi]=\inf\sigma(H_n).
\]
Smooth compactly supported antisymmetric functions form a form core:
approximate in $H^1$ and apply the finite antisymmetrization projection.
Set $\mathcal H_0=\mathcal Q_0=\mathbb C$, $q_0=0$, and $E_0=0$.
The constants in these elementary domain facts may depend on the fixed
system and on $n$.

\begin{lemma}[Monotonicity and strict-binding compactness]
\label{lem:sector-compactness}
For every fixed nuclear system, $E_{n+1}\leq E_n$ for $n\geq0$, and
$E_1<0$.  If $E_n<E_{n-1}$, there is a normalized
$\Psi\in\mathcal Q_n$ with $q_n[\Psi]=E_n$.
\end{lemma}

\begin{proof}
Approximate an $n$-electron trial state by one with compact support.
Wedge it with a normalized spin orbital in a disjoint distant ball.
Dilate the orbital to make its kinetic energy small and separate the
supports far enough to make the added repulsion small.  Disjoint spatial
supports make the energy of the normalized wedge the sum of the group
energies and their direct cross repulsion.  Dropping the added nuclear
attraction proves $E_{n+1}\leq E_n$.  A normalized orbital dilated about
one nucleus has kinetic energy $O(s^{-2})$ and attraction at most
$-c z_m s^{-1}$, proving $E_1<0$.

Suppose $\Delta=E_{n-1}-E_n>0$, and take a normalized minimizing sequence
$\psi_j$.  The form bound above gives a uniform $H^1$ bound.  Choose
nonnegative Lipschitz functions $g_S,h_S$ with
\[
 g_S^2+h_S^2=1,\quad g_S=1\text{ on }B(0,S),\quad
 g_S=0\text{ outside }B(0,2S),\quad
 |\nabla g_S|+|\nabla h_S|\leq C/S.
\]
They may be constructed as the cosine and sine of a cutoff angle.
For each subset $I$ of the particle labels, multiply $\psi_j$ by
$\prod_{i\in I}g_S(x_i)\prod_{i\notin I}h_S(x_i)$.  The squared norms of
these localized states sum to one.  The product partition identity for
the kinetic energy adds at most $Cn/S^2$ to the sum of their forms.

The all-$g_S$ component has energy at least $E_n$ times its squared
norm.  If $k<n$ coordinates carry $g_S$, slice in the remaining
coordinates.  Antisymmetry in the $k$ retained coordinates gives their
energy lower bound $E_k\geq E_{n-1}$.  Drop the other kinetic terms and
all pair terms involving an $h_S$ coordinate.  For
$S>2\max_m|R_m|$, their total nuclear attraction is bounded below by
$-2Zn/S$ times the component's squared norm.  If $p_j(S)$ is the total
squared mass of components other than all-$g_S$, it follows that
\[
 (\Delta-2Zn/S)p_j(S)
 \leq q_n[\psi_j]-E_n+Cn/S^2.
\]
The probability that any $|x_i|\geq2S$ is at most $p_j(S)$, because the
all-$g_S$ product vanishes there.  Taking first $j\to\infty$, then
$S\to\infty$, proves tightness.

Local Rellich compactness and tightness give, along a subsequence,
strong $L^2$ convergence to a normalized antisymmetric $\Psi$, weak
$H^1$ convergence, and almost-everywhere convergence.  The attraction
converges: the infinitesimal form bound applied to $\psi_j-\Psi$, with
its relative-bound parameter subsequently tending to zero, gives
\[
 \sum_i\int V(x_i)|\psi_j-\Psi|^2\longrightarrow0.
\]
Cauchy--Schwarz with this weight handles the difference of the
attractive expectations.  The kinetic energy is weakly lower
semicontinuous, and the nonnegative repulsion is lower semicontinuous
by Fatou's lemma.  Hence $\Psi$ minimizes.  Variation in the form
domain gives the weak eigen-equation at $E_n$.
\end{proof}

\subsection{A coarse number bound}

The following reciprocal-potential test is a weighted-electron argument
of the type used in charge-linear ionization bounds \citep{Lieb1984}.
We include the calculation, including the regularization needed for the
form domain.

\begin{lemma}[Reciprocal-potential bound]
\label{lem:sector-crude}
If $E_n<E_{n-1}$ and every $z_m\geq1$, then $n\leq3Z$.
\end{lemma}

\begin{proof}
Let $\Psi$ be the minimizer from Lemma~\ref{lem:sector-compactness}.
For $l,\varepsilon>0$ put
\[
 V_l(x)=\sum_m z_m(|x-R_m|^2+l^2)^{-1/2},\qquad
 U=V_l+\varepsilon,\qquad w=U^{-1}.
\]
For these fixed parameters, $w$ and its first derivatives are bounded.
Test the weak eigen-equation with
$(\sum_iw(x_i))\Psi$ and take real parts.  This summed multiplier is
symmetric; after taking this admissible test we may expand its
individual summands in the ambient form.

In the $i$th summand, the Hamiltonian on the other $n-1$ coordinates,
minus $E_n$, has nonnegative pairing after weighting by $w(x_i)$.
This follows by slicing in $x_i$ and its spin and using
$E_{n-1}\geq E_n$.  The remaining kinetic pairing is also
nonnegative.  Indeed, with $\phi=w(x_i)\Psi$,
\begin{align*}
 \operatorname{Re}\int\nabla_i\Psi\cdot
              \overline{\nabla_i(w(x_i)\Psi)}
 &=\int U(x_i)|\nabla_i\phi|^2
          -\frac12\int(\Delta V_l)(x_i)|\phi|^2\geq0,\\
 \Delta V_l(x)&=-3l^2\sum_m
              z_m(|x-R_m|^2+l^2)^{-5/2}\leq0.
\end{align*}
The integration by parts is justified by large-radius cutoffs and
$H^1$ approximation; all displayed coefficients and their needed
derivatives are bounded for fixed $l,\varepsilon$.  We obtain
\begin{equation}\label{eq:sector-weighted-pairs}
 \E_\Psi\sum_{i<j}\frac{w(x_i)+w(x_j)}{|x_i-x_j|}
 \leq\E_\Psi\sum_iw(x_i)V(x_i).
\end{equation}

Let $l\downarrow0$ at fixed $\varepsilon$.  Dominated convergence
applies, using $V/\varepsilon$ for attraction and
$2K/\varepsilon$ for a pair.  These functions have finite expectations
on the form domain.  Hence \eqref{eq:sector-weighted-pairs} holds with
$w=(V+\varepsilon)^{-1}$, and its right side is at most $n$.

For a configuration outside the null set of nuclear collisions, set
\[
 q_{im}=\frac{w(x_i)z_m}{|x_i-R_m|},\qquad
 s_i=\sum_mq_{im}=\frac{V(x_i)}{V(x_i)+\varepsilon}\leq1.
\]
The triangle inequality gives
\[
 \frac{w(x_i)+w(x_j)}{|x_i-x_j|}
 \geq\sum_m\frac{q_{im}q_{jm}}{z_m}.
\]
In fact, multiplying the right side by $|x_i-x_j|$ bounds it above by
$w(x_i)w(x_j)(V(x_i)+V(x_j))\leq w(x_i)+w(x_j)$.
Weighted Cauchy--Schwarz and $z_m\geq1$ now imply
\begin{align*}
 \sum_{i<j}\sum_m\frac{q_{im}q_{jm}}{z_m}
 &=\frac12\sum_m\frac{(\sum_iq_{im})^2-\sum_iq_{im}^2}{z_m}\\
 &\geq\frac{(\sum_i s_i)^2}{2Z}-\frac n2.
\end{align*}
Here $\sum_mq_{im}^2/z_m\leq\sum_mq_{im}^2\leq s_i^2\leq1$.
As $\varepsilon\downarrow0$, each $s_i$ tends to one almost surely,
and the squared sum is bounded by $n^2$.  Taking expectations in the
last inequality and using \eqref{eq:sector-weighted-pairs} gives
$n\geq n^2/(2Z)-n/2$, as asserted.
\end{proof}

\begin{lemma}[Reduction to a global sector minimum]
\label{lem:sector-global-minimum}
For every fixed nuclear system there is a largest integer $N$ such
that $E_N<E_{N-1}$.  It satisfies $N\leq3Z$, has a normalized
minimizer $\Psi$, and
\begin{equation}\label{eq:sector-global-minimum}
 E:=E_N=\inf_{n\geq0}E_n.
\end{equation}
Any bound on this $N$ bounds every strictly bound sector of the same
system.
\end{lemma}

\begin{proof}
The set of strict sectors is nonempty since $E_1<0$, and finite by
Lemma~\ref{lem:sector-crude}.  After its largest member there are no
strict drops.  Monotonicity therefore gives $E_n=E_N$ for $n>N$,
whereas $E_n\geq E_N$ for $n<N$.  Existence of $\Psi$ follows from
Lemma~\ref{lem:sector-compactness}.  Every other strict sector has
particle number at most $N$.
\end{proof}

For each fixed nuclear system put $E=\inf_{n\ge0}E_n$.
A \emph{$\delta$-state} is any normalized $\psi\in\mathcal Q_n$ in
a finite sector $n$ with $q_n[\psi]\le E+\delta$, where $\delta\ge0$.
Expectation in the squared law of a state includes the spin sum;
when only positions are observed, particle labels are retained unless
an unordered observation is explicitly specified.

\begin{lemma}[Energy cost of a multiplier]
\label{lem:sector-multiplier}
Let $\Psi\in\mathcal Q_n$ be a normalized eigenstate with eigenvalue $e_\Psi$.
If $F$ is a bounded real Lipschitz function of the spatial configuration,
invariant under particle permutations, then
\begin{equation}\label{eq:sector-multiplier}
 q_n[F\Psi]-e_\Psi\|F\Psi\|_2^2
 =\frac12\E_\Psi|\nabla F|^2.
\end{equation}
Here $\nabla$ denotes the full spatial gradient in $\R^{3n}$.
\end{lemma}

\begin{proof}
Both $F\Psi$ and $F^2\Psi$ belong to $\mathcal Q_n$.  Test the weak
eigen-equation with $F^2\Psi$ and expand the gradients.  Subtraction
from $q_n[F\Psi]$ cancels all multiplication potentials and
mixed-gradient terms, leaving \eqref{eq:sector-multiplier}.
\end{proof}

\subsection{Atomic sectors and neutral ground states}

In the atomic specialization $M=1$, $R_1=0$, and $z_1=Z$, we write
$E_n=E(n,Z)$. The compactness proof of Lemma~\ref{lem:sector-compactness}
applies to every normalized minimizing sequence, so the set of normalized
ground states of a strictly bound sector is compact in $L^2$.

This is Zhislin's atomic binding condition $n<Z+1$
\citep{Zhislin1960,ZhislinSigalov1965}, specialized to integer $Z$.
We give the proof needed here.

\begin{lemma}[Atomic binding]\label{lem:atomic-binding}
For integer $Z\ge1$, $E_n<E_{n-1}$ for $1\le n\le Z$.
Each of these sectors has a ground state, and every normalized minimizing
sequence has an $L^2$ convergent subsequence with normalized ground-state limit.
The multiplier identity of Lemma~\ref{lem:sector-multiplier} holds for every
such ground state with $e_\Psi=E_n$.
\end{lemma}
\begin{proof}
Strict binding follows inductively beginning with the vacuum. Suppose
$\Phi$ is a ground state in sector $n-1$. With the radial core cutoff
$g_S$ used in the compactness proof, put $F_S=\prod_i g_S(x_i)$.
Its norm tends to one, and the multiplier identity gives normalized
energy $E_{n-1}+O_Z(S^{-2})$. Wedge this state with a normalized radial
spatial orbital of fixed spin, supported in $3S<|x|<4S$, with kinetic
energy $O(S^{-2})$. The supports are separated. Newton's spherical mean
identity is
\begin{equation}\label{eq:atomic-newton}
 \frac1{4\pi}\int_{\mathbb S^2}|s\omega-x|^{-1}\,d\omega
 =\frac1{\max(s,|x|)}.
\end{equation}
It follows directly by integrating
$(s^2+|x|^2-2s|x|t)^{-1/2}/2$ over $-1<t<1$.
Every core radius is at most $2S$, so the added electron's attraction
plus direct repulsion is at most $-(Z-n+1)/(4S)$. For $n\le Z$ this
strictly dominates the kinetic and cutoff errors when $S$ is large.
Lemma~\ref{lem:sector-compactness} supplies the next ground state.
Its tightness argument gives the asserted compactness for every minimizing
sequence, completing the induction.
\end{proof}
\begin{lemma}[A minimum over all sectors]\label{lem:atomic-sector}
Every strict sector, namely every $n$ with $E_n<E_{n-1}$, satisfies
$n\leq2Z+1$.  There is consequently a largest strict index $N_*$,
and
\begin{equation}\label{a:eq:found-sector}
 Z\leq N_*\leq2Z+1\leq3Z,
 \qquad E:=\inf_{n\geq0}E_n=E_{N_*}.
\end{equation}
The infimum $E$ is attained by a normalized ground state in sector
$N_*$.
\end{lemma}

\begin{proof}
Let $\Phi$ be a ground state in a strict sector, supplied by
Lemma~\ref{lem:sector-compactness}.  For $l,\eps>0$ set
\[
 U(x)=(|x|^2+l^2)^{-1/2}+\eps,\qquad w(x)=U(x)^{-1}.
\]
For fixed parameters $w$ and its first derivatives are bounded.
Test the weak eigen-equation by $(\sum_iw(x_i))\Phi$ and take real
parts.  This is a symmetric admissible multiplier.  In its $i$th
summand, the form on the other $n-1$ coordinates is at least
$E_{n-1}\E_\Phi w(x_i)$, by slicing in $x_i$ and its spin.  The
selected-coordinate kinetic pairing is nonnegative.  Indeed, writing
$\phi=w(x_i)\Phi$ and hence $\Phi=U(x_i)\phi$, integration by parts
gives
\begin{align*}
 \operatorname{Re}\int\nabla_i\Phi\cdot
                    \overline{\nabla_i(w(x_i)\Phi)}
 &=\int U(x_i)|\nabla_i\phi|^2
       -\frac12\int\Delta U(x_i)|\phi|^2\geq0,\\
 \Delta U(x)&=-3l^2(|x|^2+l^2)^{-5/2}\leq0.
\end{align*}
The identity follows first for smooth compactly supported functions;
the bounded coefficients for fixed $l,\eps$, approximation in $H^1$,
and large-radius cutoffs give it in the present domain.
Since $E_{n-1}-E_n>0$, discarding that additional nonnegative term
leaves
\begin{equation}\label{a:eq:found-weighted}
 \E_\Phi\sum_{i<j}\frac{w(x_i)+w(x_j)}{|x_i-x_j|}
 \leq \E_\Phi\sum_i w(x_i)\frac Z{|x_i|}.
\end{equation}
Let $l\downarrow0$ at fixed $\eps$.  Domination by
$2K(x_i-x_j)/\eps$ for pairs and $V(x_i)/\eps$ for attraction,
which are integrable on the form domain, justifies the limit.
Then $w(x)=|x|/(1+\eps|x|)$.  As $\eps\downarrow0$, both sides
increase and the right side tends to $nZ$.  The left side tends to
the expectation of
$\sum_{i<j}(|x_i|+|x_j|)/|x_i-x_j|$, which is at least
$n(n-1)/2$ by the triangle inequality.  Thus $n\leq2Z+1$.

The strict indices form a nonempty finite set by
Lemma~\ref{lem:atomic-binding}.  After its largest element $N_*$
there can be no further strict drop; monotonicity implies
$E_n=E_{N_*}$ for every $n>N_*$.  For $n<N_*$ monotonicity gives
$E_n\geq E_{N_*}$.  This proves \eqref{a:eq:found-sector}, without
any assumption on possible plateaus before $N_*$.  Attainment
follows once more from Lemma~\ref{lem:sector-compactness}.
\end{proof}

